%% CV for Aalborg Portland - Group Electrical \& Automation Engineer (Dansk)
%% Compiled with: lualatex

\documentclass[11pt,a4paper,sans]{moderncv}
\moderncvstyle{banking}
\moderncvcolor{blue}

\renewcommand*{\firstnamestyle}[1]{{\fontsize{34}{36}\bfseries\upshape\color{color1}#1}}
\renewcommand*{\lastnamestyle}[1]{{\fontsize{34}{36}\bfseries\upshape\color{color1}#1}}
\renewcommand*{\sectionstyle}[1]{{\sectionfont\color{color1}#1}}

\usepackage[utf8]{inputenc}
\usepackage{hyperref}
\hypersetup{
    colorlinks=true,
    linkcolor=blue,
    filecolor=magenta,
    urlcolor=blue,
    pdftitle={Jacob El-Omar - CV},
    pdfpagemode=FullScreen,
}
\usepackage[scale=0.80]{geometry}
\usepackage{import}

\name{Jacob}{El-Omar}
\address{Aalborg, Danmark}{}{}
\phone[mobile]{+45 53 35 27 82}
\email{elomarjc@gmail.com}
\extrainfo{\href{https://linkedin.com/in/jacob-el-omar}{LinkedIn} \textbullet{} \href{https://github.com/elomarjc}{GitHub}}

\begin{document}

\makecvtitle

% ============================================================
%     PROFIL
% ============================================================

\vspace{6pt}
\small{Nyuddannet civilingeniør (cand.polyt.) i elektroniske systemer fra Aalborg Universitet med speciale i reguleringsteknik, modellering og industriel automation. Erfaren i regulator-design (LQR, PID), systemidentifikation og sensorintegration samt praktisk fremstilling (CNC og printudlægning). En analytisk og samarbejdsorienteret ingeniør trænet i problembaseret læring (PBL), klar til at understøtte jeres el- og automationsprojekter i gruppen.}

% ============================================================
%     FAGLIGE KOMPETENCER
% ============================================================

\section{Faglige Kompetencer}
\vspace{1pt}
\begin{itemize}

\item \textbf{Reguleringsteknik \& Modellering}: Design af reguleringssløjfer, LQR, systemidentifikation, modellering af fysiske systemer i MATLAB og Simulink.

\item \textbf{Automation \& Industriel IT}: Kendskab til PLC-arkitekturer, ladder og struktureret tekst programmering, samt kommunikationsprotokoller (såsom OPC-UA).

\item \textbf{Elektronik \& Hardware}: Printudlægning (Altium Designer), sensorintegration (I2C, SPI, analoge sensorer), kalibrering samt praktisk mekanisk montage (CNC).

\item \textbf{Signalbehandling \& IoT}: Python-dataanalyse, adaptive støjfiltreringsalgoritmer (LMS, RLS) samt kendskab til trådløse teknologier (UWB).

\item \textbf{Projektstyring \& Kommunikation}: Erfaren med agilt tværfagligt samarbejde og systematisk fejlsøgning på komplekse el-tekniske anlæg.

\end{itemize}

% ============================================================
%     ERHVERVSERFARING
% ============================================================

\section{Erhvervserfaring}
\vspace{3pt}
\begin{itemize}

\item{\cventry{Sep 2024--Jun 2025}{AI/ML Engineer (Praktik)}{Ai Health Highway}{Remote}{}{\vspace{1pt}
\begin{itemize}
    \item Udviklede og testede algoritmer i Python til realtids-signalbehandling i smarte stetoskoper.
    \item Kalibrerede støjreduktionsmodeller og validerede signalintegriteten mod systemkrav.
    \item Samarbejdede tæt med hardware- og softwareudviklere om integration af sensordata.
\end{itemize}}}

\vspace{3pt}

\item{\cventry{Jan 2026--Nu}{IT- \& Webudvikler (Frivillig)}{Dawah Danmark}{Aalborg, Danmark}{}{\vspace{1pt}
\begin{itemize}
    \item Ansvarlig for serverdrift, databaseadministration og integrationer af webapplikationer.
\end{itemize}}}

\vspace{3pt}

\item{\cventry{Jun 2023--Aug 2024}{Lager- og Logistikmedarbejder}{Lyn Express}{Aalborg, Danmark}{}{\vspace{1pt}
\begin{itemize}
    \item Styrede sortering og logistik under højt tidspres i et fysisk krævende lager- og logistikmiljø.
\end{itemize}}}

\end{itemize}

% ============================================================
%     UDDANNELSE
% ============================================================

\section{Uddannelse}
\vspace{1pt}
\begin{itemize}

\item{\cventry{Sep 2023--Jun 2025}{Civilingeniør, cand.polyt. (Elektroniske Systemer)}{Aalborg Universitet}{Aalborg, Denmark}{}{\vspace{1pt}
Speciale: ``Adaptive Noise Cancellation For Electronic Stethoscopes.'' Støjfiltrering og realtids-signalbehandling på indlejrede systemer.
}}

\vspace{3pt}

\item{\cventry{Sep 2020--Aug 2023}{Bachelor i teknisk videnskab (Elektronik og IT)}{Aalborg University}{Aalborg, Danmark}{}{\vspace{1pt}
Specialisering i reguleringsteknik og automation. Afsluttede projekter omfattede:
\begin{itemize}
    \item \textbf{Bachelorprojekt}: Modelopbygning og regulering af en traverskran (Gantry Crane).
    \item \textbf{5. semester}: Kollisionsforebyggende navigationssystem til mobile robotter (MiR200) baseret på UWB og regressionsalgoritmer.
    \item \textbf{4. semester}: Indlejret faldsikringssystem på Cypress PSoC i C.
\end{itemize}
}}

\end{itemize}

% ============================================================
%     UDVALGTE PROJEKTER
% ============================================================

\section{Udvalgte Projekter}
\vspace{1pt}
\begin{itemize}

\item{\textbf{Regulering af Traverskran (Bachelorprojekt, Universitetsprojekt)}: Udviklede matematisk model og reguleringssløjfer (PID/LQR) for automatiseret positionering og svingningsdæmpning af en fysisk modeltraverskran (Gantry Crane) med sensorfeedback.}

\item{\textbf{AAUSAT6 CubeSat ADCS Design (Universitetsprojekt)}: Udviklede simulationsmodeller og attitude-regulatorer (LQR, robuste B-dot algoritmer) til CubeSat-satellitstabilisering i MATLAB/Simulink.}

\item{\textbf{Center of Mass Kalibrering (Universitetsprojekt, 3-DoF Simulator)}: Udviklede estimationsalgoritmer til Center of Mass ved brug af kinematik og Kalman-filtrering. Stod for PCB-layout i Altium Designer og mekanisk CNC-fræsning.}

\end{itemize}

% ============================================================
%     ANERKENDELSER
% ============================================================

\section{Anerkendelser}
\vspace{1pt}
\begin{itemize}
\item{\textbf{1. plads, RoboCup-turneringen} -- Aalborg Universitet (2020). Udviklede navigations- og reguleringsalgoritmer til en mobil robot.}
\end{itemize}

% ============================================================
%     SPROG
% ============================================================

\section{Sprog}
\vspace{1pt}
\begin{itemize}
\item Dansk (modersmål), Engelsk (flydende).
\end{itemize}

% ============================================================
%     REFERENCER
% ============================================================

\section{References}
\vspace{1pt}
\begin{itemize}
\item{Afleveres efter aftale.}
\end{itemize}

\end{document}
